\Needspace{13\baselineskip}
\subsection{Varying divisorial coefficients}

For a fixed support, the residual equation depends smoothly on the coefficients as long as the residual background stays in $\mathscr P$. The smooth remainders then give convergence in capacity of the solutions, and they determine all generic Lelong numbers of the limit on every modification. Keep $D_j$, $s_j$ and $\eta_j$ as above, and put $\sigma_j=\log|s_j|^2$.

For $a=(a_1,\ldots,a_N)\in\mathbb R_{\geq0}^N$ put
\[
 D_a=\sum_j a_jD_j,\qquad
 \sigma_a=\sum_{a_j>0}a_j\sigma_j,\qquad
 \beta_a=\theta-\sum_j a_j\eta_j,
\]
where zero coefficients are omitted from the logarithmic sum. Let $\mathcal A$ be the set of such $a$ for which $\beta_a$ has a smooth positive representative.

\Needspace{14\baselineskip}
\begin{proposition}\label{prop:finite-divisorial-equation-limits}
There are unique smooth functions $v_a$, $a\in\mathcal A$, such that
\[
 \int_Xv_a\,\mathrm{d}\mu=0,\qquad
 \Omega_a=\beta_a+\ddc v_a>0,\qquad
 \Omega_a^n=c(a)\mu,
 \quad c(a)=\mathcal V_\mu(\beta_a)>0.
\]
The maps $a\mapsto v_a$ and $a\mapsto c(a)$ are smooth, and every smooth norm of $v_a$ is bounded on compact subsets of $\mathcal A$. Here smoothness at a parameter with a zero coefficient means that the family is the restriction of a smooth family defined on an open subset of $\mathbb R^N$.

Set $u_a=\sigma_a+v_a$. If $a^{(k)}\in\mathcal A$ and $a^{(k)}\to a\in\mathcal A$, then $u_{a^{(k)}}\to u_a$ in $L^1$ and in capacity, and smoothly on compact subsets of $X\setminus\bigcup_jD_j$. Moreover,
\[
 \bigl\|\mathrm{MA}^{\mathrm{ext}}_\theta(u_{a^{(k)}})
       -\mathrm{MA}^{\mathrm{ext}}_\theta(u_a)\bigr\|_{\mathrm{TV}}
       =|c(a^{(k)})-c(a)|\longrightarrow0,
\]
where $\|\cdot\|_{\mathrm{TV}}$ is the total variation norm. For every modification $\pi:Y\to X$ with $Y$ smooth and every prime divisor $E\subset Y$,
\[
 \nu(u_a\circ\pi,E)=\sum_j a_j\operatorname{ord}_E(\pi^*D_j),
\]
where $\operatorname{ord}_E(\pi^*D_j)$ is the vanishing order along $E$ of the pullback of a local equation of $D_j$. In particular, the limit $u_a$ has exactly the ordinary singularity type of $\sigma_a$, and no further singularity appears in this finite family.
\end{proposition}
\begin{proof}
The set of parameters in $\mathbb R^N$ for which $\beta_a$ has a smooth positive representative is open, and the smooth equation gives $v_a$ and $c(a)$ on it. The argument in the proof of \cref{prop:gaud-volume-first-variation} gives smooth dependence of both quantities on the parameter, with the fixed mean normalization; uniqueness makes the local families agree. Covering a compact set of parameters by finitely many such neighborhoods gives the asserted bounds.

The Poincar\'e--Lelong formula gives
\[
 \theta+\ddc u_a=[D_a]+\Omega_a,
\]
so $u_a\in\PSH(X,\theta)$. Its extended measure equals $\Omega_a^n=c(a)\mu$ off the union of the divisors, and both measures vanish on that union; this also covers the case $D_a=0$. Each fixed $\sigma_j$ is locally integrable, so the convergence of the coefficients and the smooth convergence of $v_a$ give convergence in $L^1$. Off the union of the divisors every $\sigma_j$ is smooth, which gives smooth convergence there. For convergence in capacity, remove a neighborhood of this analytic set of arbitrarily small Hermitian capacity; on its compact complement the convergence is uniform. Since the union of the divisors has capacity zero, this proves convergence in capacity. The identity for the total variation distances follows from the equations and $\mu(X)=1$.

Finally, after pullback the logarithmic potential $\sigma_j\circ\pi$ has generic Lelong number $\operatorname{ord}_E(\pi^*D_j)$ along $E$, whereas $v_a\circ\pi$ is smooth. Additivity of Lelong numbers for the finite sum proves the formula on every modification. Since $v_a$ is smooth also at the limiting parameter, $u_a$ has the ordinary singularity type of $\sigma_a$.
\end{proof}

The smooth bounds above use that the limiting parameter lies in $\mathcal A$. When the sections and weights move, \cref{thm:moving-divisorial-potential-stability} gives $L^1$ convergence of the potentials normalized by $\sup_X=0$; when the residual background is fixed, it also allows uniformly bounded measurable right-hand sides converging in $L^1$. On a modification, the pullback of a positive probability density can vanish along the exceptional divisor, so the smooth equation with positive density no longer applies directly. For a minimal exceptional pole, \cref{prop:exceptional-density-obstruction} gives a vanishing order of the density that is necessary for a smooth residual, but this order alone does not ensure smoothness.

Multiplier ideals need not stabilize in this family. At a smooth point of $D_j$ outside the other divisors, take a local coordinate $z_1$ defining $D_j$. There $u_a=a_j\log|z_1|^2+O(1)$, so
\[
 \mathcal I(u_a)=(z_1^{\lfloor a_j\rfloor})
\]
by the one-variable integral in \cref{ex:ideal-disc} and Fubini's theorem. Thus, if the parameter stays in $\mathcal A$ while $a_j$ increases to an integer from below, this local ideal jumps at the limit. Convergence of the Lelong numbers to their limiting values does not imply that the multiplier ideals are eventually equal to those of the limit.
